\section{Matter and the Fields (P4, P14, P20, P22, P23, P25)}
\label{sec:matter}

\subsection{The monist spinor and spin from topology (P4)}

Matter is built from the same tones, with no new primitive. A spinor is a graded
complex of ternary vibes (vertices, edges as relational vibes, faces), not a richer
tone on one vibe, with its complex phase emergent from a finite clock-tone. The
Kahler-Dirac zero-mode count equals the surface Betti sum exactly: $1$ for a disk, $2$
for a cylinder, $4$ for a torus. \textbf{Where we landed:} matter is monist and spin is
topological. The chirality wall (Nielsen-Ninomiya) is threaded: the overlap operator
gives a single fermion species with exact lattice chiral symmetry (Ginsparg-Wilson
residual $6\times10^{-16}$, machine zero), where the naive operator has four doublers.
This is at the effective (Dirac-on-a-complex) level, the bottom-up ternary-cluster
construction being the open frontier.

\subsection{Mass and the relativistic dispersion (P14)}

A mass term in the lattice Dirac Hamiltonian opens a spectral gap equal to the rest
energy and bends the dispersion into the relativistic form. \textbf{Where we landed:}
the gap equals the mass exactly, the dispersion is $\omega^2 = k^2 + m^2$ (slope near
one, intercept $m^2$), and at $m=0$ the gap vanishes and $\omega = |k|$, a massless
light-cone mode.

\subsection{The photon (P20) and the gauge operator from the action (P23)}

The free U(1) gauge field is a proper photon: on a periodic lattice about one third of
its modes are exact gauge zero modes (the gauge freedom) and the rest are the two
transverse polarizations, with a massless spectrum (the smallest physical $\omega^2$
shrinks as the lattice grows) that a mass term would gap. \textbf{Where we landed:} the
photon operator is not put in by hand. It is the small-field limit of the Wilson gauge
action, $1 - \cos\theta \to \tfrac12\theta^2$: the ratio of the Wilson action to the
Maxwell action converges to one as the field shrinks ($0.926, 0.988, 0.997, 0.9999$).
So the gauge kinetic operator is derived from the discrete action.

\subsection{The Higgs (P22) and electroweak breaking (P25)}

Mass is generated, not inserted. A scalar with a Mexican-hat potential settles at a
nonzero vacuum value $v$ (the symmetric phase has $v=0$, the broken phase $v=1$), and a
gauge field coupled to it eats the Goldstone mode and becomes massive: the photon gap
goes from $0$ to exactly $(gv)^2$. \textbf{Where we landed:} this extends to the
specific Standard-Model pattern. A Higgs doublet breaks $SU(2)\times U(1)_Y$ to
$U(1)_{\mathrm{EM}}$, giving the W and Z mass and leaving the photon massless, with the
Weinberg relation $m_W = m_Z\cos\theta_W$. With the measured couplings it reproduces the
observed masses: $m_W = 80.0$ GeV, $m_Z = 91.3$ GeV, $\sin^2\theta_W = 0.233$. So scalar,
spinor, and vector fields are all present, with mass from symmetry breaking. The rest of
the Standard-Model content (three generations, the Yukawa spectrum, the gauge group
itself) is the larger program.
